Let \(d=p(p-1)/2\), and let \(g_r\) denote the map from a raw moment vector to the \(d\)-vector of upper-triangular entries of the recovered signed jump \(Z^\star=K^\star-(K^\star)^{\mathsf T}\), using the rank-\(r\) truncated image construction.  On the valid event \(\mathcal V_n\), the definitions give, for every \(i<j\),
\[
 \lambda_j^\star A_{ij}^\star-\lambda_i^\star A_{ji}^\star
 =\max\{Z_{ij}^\star,0\}-\max\{-Z_{ij}^\star,0\}
 =Z_{ij}^\star.
\tag{13}
\]
The additional invertibility condition in \(\mathcal V_n\), together with
\(\mu^\star=(I_p-A^\star/\beta^\star)\lambda^\star\), shows that the intensity computed inside \(\zeta(\widehat\theta_n)\) is exactly \(\lambda^\star\).  On \(\mathcal V_n^c\), both \(\widehat z\) and \(\zeta(1,0,\mathbf1_p)\) equal zero.  Therefore the totalized definitions satisfy the exact identity
\[
 \widehat z=g_{\widehat r_n}(\widehat\xi_n)=\zeta(\widehat\theta_n)
\quad\text{on }\mathcal V_n,
\quad
 \widehat z=\zeta(\widehat\theta_n)=0
\quad\text{on }\mathcal V_n^c.
\tag{14}
\]

The projector calculation (4)--(6) in the proof of \(\text{thm:rank-adaptive-estimation}\) shows that the truncated active image has a differentiable local frame whenever its singular-value gap is positive.  Conditional on such a frame, \(F,q,\beta^\star,E,L^\star,H^\star,K^\star\), and \(Z^\star\) are obtained by matrix multiplication, inversion at nonsingular matrices, the scalar logarithm on \((0,1)\), and a finite matrix polynomial.  Hence \(g_r\) is Fr\'{e}chet differentiable at every population point in \(\mathcal M_{\Delta,p}(\kappa)\).  Equation (13) proves that this derivative is the derivative of the composite \(\zeta\circ\Phi_{\Delta,r}\circ\mathcal C\), even when \(Z_{ij}=0\).  In particular, differentiability of either positive part separately is neither true nor needed.

All derivative bounds used below are uniform.  Indeed, equations (11)--(12), the decay margin, the intensity gap, and the parameter envelope give common bounds for every inverse and divisor.  The eighth-moment envelope bounds the input moments, and the resolvent identity (5) bounds the derivative of the active projector and its local modulus.  Since \(r\in\{1,\ldots,p\}\), taking the maximum over the finitely many rank strata produces a common derivative bound and a common modulus of continuity.  Taylor's formula therefore gives, on the uniformly likely event \(\{\widehat r_n=r(\theta)\}\cap\mathcal V_n\),
\[
 \sqrt n\{\widehat z-\zeta(\theta)\}
 =J_z(\theta)\sqrt n\{\widehat\xi_n-\xi(\theta)\}+o_P(1),
\qquad
 J_z(\theta)=D(\zeta\circ\Phi_{\Delta,r(\theta)}\circ\mathcal C)_{\xi(\theta)},
\tag{15}
\]
where the remainder is uniform over the margin class.  The uniform rank and validity probabilities follow from \(\text{thm:rank-adaptive-estimation}\).  Applying the bounded linear maps \(J_z(\theta)\) to the uniform Gaussian approximation in \(\text{lem:bin-moment-limit}\), and then using (14)--(15), proves
\[
 \sup_{\theta:P_\theta\in\mathcal M_{\Delta,p}(\kappa)}
 d_{\mathrm{BL}}\!\left(
 \mathcal L_\theta[\sqrt n\{\zeta(\widehat\theta_n)-\zeta(\theta)\}],
 \mathcal N\{0,W(\theta)\}\right)\longrightarrow0,
\]
with \(W(\theta)=J_z(\theta)\Omega(\theta)J_z(\theta)^{\mathsf T}\), as stated.

We next verify both finite-sample positive semidefiniteness and consistency of the covariance used for calibration.  Put \(N_n=n-2\) and \(v_s=Y_s-\overline Y_n\).  The Bartlett definition can be rewritten as
\[
 \widehat\Omega_n
 =\frac1{N_n}\sum_{s,t=0}^{N_n-1}
 \left(1-\frac{|s-t|}{b_n}\right)_{+}v_sv_t^{\mathsf T}.
\tag{16}
\]
For integer \(b_n\geq1\), the scalar matrix in (16) is positive semidefinite because
\[
 \left(1-\frac{|s-t|}{b_n}\right)_{+}
 =\frac1{b_n}\sum_{\ell\in\mathbb Z}
 \mathbf1_{\{s\leq\ell<s+b_n\}}
 \mathbf1_{\{t\leq\ell<t+b_n\}},
\tag{17}
\]
which is a Gram representation.  Thus \(a^{\mathsf T}\widehat\Omega_n a\geq0\) for every \(a\), and the sandwich defining \(\widehat W_n\) is positive semidefinite on \(\mathcal V_n\); on \(\mathcal V_n^c\) it is the zero matrix and is again positive semidefinite.  The Bartlett kernel is even, Lipschitz, supported on \([-1,1]\), and equals one at zero, so \(\text{lem:bin-moment-limit}\) yields uniform consistency of \(\widehat\Omega_n\).  Derivative continuity, (15), and the uniform probability of \(\mathcal V_n\) now give
\[
 \sup_{\theta:P_\theta\in\mathcal M_{\Delta,p}(\kappa)}
 P_\theta\{\lVert\widehat W_n-W(\theta)\rVert_{\mathrm{op}}>\varepsilon\}
 \longrightarrow0
\tag{18}
\]
for every \(\varepsilon>0\).

For completeness, we verify that singular covariance matrices do not invalidate Gaussian-maximum calibration.  The covariance family has a uniformly nondegenerate coordinate.  Define the raw-moment direction
\[
 v_\theta=\{0,\operatorname{vec}(\Gamma_1),\operatorname{vec}(\Gamma_2)\}.
\]
Along \(\xi+s v_\theta\), the mean is unchanged and both centered lag covariances are multiplied by \(1+s\).  Consequently the active image is unchanged, \(R\) is multiplied by \(1+s\), and
\[
 F(s)=O^{\mathsf T}(1+s)\Gamma_2\{(1+s)R\}^{\mathsf T}
 \left[\{(1+s)R\}\{(1+s)R\}^{\mathsf T}\right]^{-1}=F.
\]
Thus \(q,E,L^\star,H^\star\) are unchanged, whereas \(K^\star,Z^\star\), and \(g_r\) are multiplied by \(1+s\).  Hence
\[
 J_z(\theta)v_\theta=\zeta(\theta).
\tag{19}
\]
The eighth-moment bound gives \(\lVert v_\theta\rVert_2\leq C_0\) uniformly.  Because the support is nonempty and \(\text{ass:edge-gap}\) holds, some coordinate \(a=a(\theta)\) satisfies \(|\zeta_a(\theta)|\geq\kappa\).  Equation (19) implies that the corresponding row \(j_a\) of \(J_z(\theta)\) obeys \(\lVert j_a\rVert_2\geq\kappa/C_0\).  Assumption \(\text{ass:variance-gap}\) then gives
\[
 W_{aa}(\theta)=j_a\Omega(\theta)j_a^{\mathsf T}
 \geq\kappa\lVert j_a\rVert_2^2
 \geq c_0:=\kappa^3/C_0^2>0.
\tag{20}
\]
The uniform derivative and long-run-variance bounds also give \(\lVert W(\theta)\rVert_{\mathrm{op}}\leq C_1\).

Let \(G\sim\mathcal N(0,W(\theta))\) and \(M=\max_{1\leq a\leq d}|G_a|\).  Equations (20) and the upper bound imply the following uniform anti-concentration estimate.  For \(0<\varepsilon<1\) and \(x\geq0\), if \(x\leq\sqrt\varepsilon\), choose the coordinate in (20) to obtain
\[
 P(x<M\leq x+\varepsilon)
 \leq P\{|G_a|\leq\sqrt\varepsilon+\varepsilon\}
 \leq C_2\sqrt\varepsilon.
\]
If \(x>\sqrt\varepsilon\), a union bound applies.  A centered normal coordinate with variance \(s^2>0\) has absolute-value density at \(y>0\) bounded, uniformly in \(s\), by \(C/y\); a zero-variance coordinate contributes nothing.  Therefore
\[
 P(x<M\leq x+\varepsilon)
 \leq\sum_{a=1}^dP(x<|G_a|\leq x+\varepsilon)
 \leq C_3d\frac{\varepsilon}{x}
 \leq C_3d\sqrt\varepsilon.
\tag{21}
\]
This proves a uniform modulus of continuity for the Gaussian-maximum distribution, including singular \(W(\theta)\).

By (18), on \(\mathcal V_n\) the covariance \(\widehat W_n\) lies with probability tending uniformly to one in an arbitrarily small neighborhood of \(W(\theta)\).  The spectral decomposition shows that the PSD square-root map is continuous; on the compact set \(\{W\succeq0:\lVert W\rVert_{\mathrm{op}}\leq C_1+1\}\) it is uniformly continuous.  Coupling \(\widehat W_n^{1/2}U\) and \(W(\theta)^{1/2}U\) with one \(U\sim\mathcal N(0,I_d)\), and using (21), shows that the conditional distribution function of the plug-in Gaussian maximum converges uniformly to the target distribution function.  Hence its conditional \((1-\alpha)\)-quantile \(c_{n,1-\alpha}\) satisfies the size inequality
\[
 \inf_{\theta:P_\theta\in\mathcal M_{\Delta,p}(\kappa)}
 P_\theta\!\left[
 \max_{i<j}\left|\sqrt n\{\widehat z_{ij}-\zeta_{ij}(\theta)\}\right|
 \leq c_{n,1-\alpha},\ \mathcal V_n\right]
 \geq1-\alpha-o(1)-\sup_\theta P_\theta(\mathcal V_n^c).
\tag{22}
\]
Here the passage from bounded-Lipschitz convergence to the indicator in (22) is justified by upper and lower Lipschitz approximations to the rectangle, with the boundary probability controlled by (21).  On \(\mathcal V_n^c\), the corrected definition sets every \(I_{ij,n}=\mathbb R\), so simultaneous coordinate coverage holds automatically.  Combining this fallback with (22) and \(\sup_\theta P_\theta(\mathcal V_n^c)\to0\) proves
\[
 \liminf_{n\to\infty}\inf_{\theta:P_\theta\in\mathcal M_{\Delta,p}(\kappa)}
 P_\theta\{\zeta_{ij}(\theta)\in I_{ij,n}\text{ for all }i<j\}
 \geq1-\alpha.
\]

Finally, acyclicity and nonnegative excitation imply that for each unordered pair at most one of \(A_{ij},A_{ji}\) is positive, while \(\lambda_i,\lambda_j>0\).  Thus \(\zeta_{ij}>0\) means the true edge is \(j\to i\), \(\zeta_{ij}<0\) means it is \(i\to j\), and \(\zeta_{ij}=0\) means neither edge is present.  On the simultaneous interval event, every true pairwise choice is admitted by the sign-inversion rules; the true graph is acyclic and hence lies in \(\mathcal C_n(1-\alpha)\).  On \(\mathcal V_n^c\), the graph set is all labeled acyclic graphs, so the same inclusion is automatic.  Therefore
\[
 \liminf_{n\to\infty}\inf_{\theta:P_\theta\in\mathcal M_{\Delta,p}(\kappa)}
 P_\theta\{G(A)\in\mathcal C_n(1-\alpha)\}\geq1-\alpha.
\]
The optional full-parameter Wald object is not used anywhere in this argument; its full-space fallback is necessary precisely where the positive-part map has no Fr\'{e}chet derivative.